\documentclass[11pt]{article}
\usepackage[margin=1in]{geometry}
\usepackage{amsmath,amssymb,amsthm,booktabs}
\usepackage[hidelinks]{hyperref}
\usepackage{microtype}

\newtheorem{theorem}{Theorem}
\newtheorem{lemma}[theorem]{Lemma}
\newtheorem{proposition}[theorem]{Proposition}
\newtheorem{corollary}[theorem]{Corollary}
\theoremstyle{definition}
\newtheorem{definition}[theorem]{Definition}
\newtheorem{remark}[theorem]{Remark}

\title{A Density-One Composite-Modulus Backbone in the Exact Type-II Erd\H{o}s--Straus Depth Spectrum}
\author{Skip Potter}
\date{Draft -- August 2026}

\begin{document}
\maketitle

\begin{abstract}
For a prime $p\equiv1\pmod4$, define
\[
\delta_{\rm sq}(p)=\min\{A\ge1:\exists D\mid A^2,\ 4A-1\mid p+4D\},
\]
the exact first-hit depth in the classical square-divisor coordinate for standard prime Type-II Erd\H{o}s--Straus solutions.

Let $q\equiv1\pmod4$ be prime and put
\[
A_q=\frac{3q+1}{4}.
\]
We prove that if $A_q$ has a prime divisor congruent to $2\pmod3$, then the deliberately composite target modulus $4A_q-1=3q$ supports infinitely many primes $p$ with
\[
\delta_{\rm sq}(p)=A_q.
\]
More generally any $D\mid A_q^2$ with $D\equiv2\pmod3$ gives an infinite exact-depth prime progression.

We then prove that this factor condition holds for relative density $1$ of primes $q\equiv1\pmod4$.  Consequently, if $N_{\rm comp}(X)$ counts depths $A\le X$ with composite $4A-1$ that are realized as exact square-completion depth for infinitely many primes, then
\[
\boxed{N_{\rm comp}(X)\ge\left(\frac23+o(1)\right)\frac{X}{\log X}.}
\]
The proof is elementary apart from standard prime-distribution theorems: it uses a universal safe residue, CRT, Dirichlet's theorem, and the divergence of reciprocal primes in the class $2\pmod3$.

We also describe a finite CRT certificate method for arbitrary composite target layers and record independently verifiable full-layer examples at $A=649$ and $A=22339$.  These results concern Type-II exact depth only; they do not prove universal Type-II solvability or the Erd\H{o}s--Straus conjecture.
\end{abstract}

\section{Introduction}
The Erd\H{o}s--Straus conjecture asks whether
\[
\frac4n=\frac1x+\frac1y+\frac1z
\]
has a solution in positive integers for every $n\ge2$.  For prime denominators, standard solutions split into Type I and Type II.  We study the exact first-hit spectrum of Type II in the classical divisor-square coordinate
\[
D\mid A^2,\qquad 4A-1\mid p+4D. \tag{1}
\]
The mechanism (1) is prior art in the Th\'epault--Mizony lineage and is not claimed here.

For primes $p\equiv1\pmod4$, define
\[
\delta_{\rm sq}(p)=\min\{A\ge1:\exists D\mid A^2\text{ satisfying (1)}\}. \tag{2}
\]
Recent public work of FCF/CENTL develops an exact-depth spectrum for this square-completed Type-II route, including an infinite prime-modulus backbone, structural gaps, and a composite residual core.  The present note gives an elementary backbone lying entirely in the composite target regime.

The target moduli in our family are
\[
4A_q-1=3q,
\qquad
A_q=\frac{3q+1}{4},
\qquad q\equiv1\pmod4\text{ prime}.
\]
The factor $3$ is decisive: it is the smallest positive integer congruent to $3\pmod4$, so the new prime factor $q$ cannot occur in any earlier modulus $4B-1$.  A target divisor $D\equiv2\pmod3$ then makes the target residue agree, on every possible common factor, with one universal safe residue at all earlier layers.

The density theorem is also elementary.  For almost every prime $q\equiv1\pmod4$, the integer $(3q+1)/4$ has a prime divisor congruent to $2\pmod3$.  This follows from finite CRT equidistribution and the divergence of the reciprocal sum over primes $2\pmod3$, without a uniform sieve estimate.

\section{The square-divisor coordinate}
We use the standard completeness fact.

\begin{proposition}[complete prime Type-II coordinate]\label{prop:complete}
For a prime $p\equiv1\pmod4$, a standard Type-II Erd\H{o}s--Straus solution exists if and only if there are positive integers $A,D$ such that
\[
D\mid A^2,\qquad p\nmid D,\qquad 4A-1\mid p+4D.
\]
\end{proposition}

\begin{remark}
Proposition~\ref{prop:complete} is prior mathematics.  It is included only to make clear that $\delta_{\rm sq}$ measures the complete standard prime Type-II family, not a sufficient subfamily.  The parameter $A$ is a square-completion depth coordinate and should not be conflated with other search-depth coordinates in the literature.
\end{remark}

\section{A universal safe residue}
\begin{lemma}[universal safe residue]\label{lem:safe}
For every $B\ge1$ and every divisor $E\mid B^2$,
\[
4B-1\nmid4E+1.
\]
Equivalently, residue $1$ modulo $4B-1$ is never a square-completed Type-II trap residue at depth $B$.
\end{lemma}

\begin{proof}
Let
\[
g=\gcd(B,E),\qquad B=ga,\qquad E=gd,\qquad\gcd(a,d)=1.
\]
Since $E\mid B^2$, we have $d\mid g$; write $g=dh$.  Thus
\[
B=dha,\qquad E=d^2h.
\]
If $4B-1\mid4E+1$, then modulo $4B-1$,
\[
E\equiv-B,
\]
so
\[
4B-1\mid E+B=dh(d+a).
\]
Because every prime divisor of $dh$ divides $B$,
\[
\gcd(4B-1,dh)=1.
\]
Hence $4B-1\mid d+a$, impossible since
\[
4B-1=4dha-1>d+a.
\]
\end{proof}

\section{A structural composite exact-depth theorem}
\begin{theorem}\label{thm:structural}
Let $q\equiv1\pmod4$ be prime and put
\[
A_q=\frac{3q+1}{4}.
\]
If there exists
\[
D\mid A_q^2,
\qquad
D\equiv2\pmod3,
\]
then $4A_q-1=3q$ is composite and infinitely many primes $p\equiv1\pmod4$ satisfy
\[
\boxed{\delta_{\rm sq}(p)=A_q.}
\]
\end{theorem}

\begin{proof}
Set
\[
m=3q=4A_q-1,
\qquad
a\equiv-4D\pmod m.
\]
Since $D\equiv2\pmod3$,
\[
a\equiv1\pmod3. \tag{3}
\]
Every integer $n\equiv a\pmod m$ has a target witness at depth $A_q$.

Let $B<A_q$ and $m_B=4B-1$.  If $q\mid m_B$, write $m_B=qs$.  Because $q\equiv1\pmod4$ and $m_B\equiv3\pmod4$, one has $s\equiv3\pmod4$ and hence $s\ge3$.  Therefore
\[
m_B\ge3q=m,
\]
contradicting $B<A_q$.  Thus
\[
q\nmid m_B,
\qquad
\gcd(m,m_B)\mid3 \tag{4}
\]
for every earlier depth.

Let
\[
L=\operatorname{lcm}_{1\le B<A_q}(4B-1).
\]
By (3)--(4), the system
\[
n\equiv-4D\pmod{3q},
\qquad
n\equiv1\pmod L,
\qquad
n\equiv1\pmod4 \tag{5}
\]
is CRT-compatible.  Every member of the resulting progression is congruent to $1$ modulo every earlier $4B-1$, so Lemma~\ref{lem:safe} excludes all earlier Type-II witnesses.  The target witness exists by construction.

The progression is reduced.  Every prime factor inherited from $L$ sees residue $1$.  The only new odd prime factor is $q$, and $q\nmid D$ because $q\nmid A_q$.  The residue modulo $4$ is $1$.  Dirichlet's theorem therefore supplies infinitely many prime members of (5), each having exact depth $A_q$.
\end{proof}

\begin{corollary}[prime-factor criterion]\label{cor:pf}
If $A_q$ has a prime divisor
\[
\ell\equiv2\pmod3,
\]
then $A_q$ is realized as $\delta_{\rm sq}(p)$ by infinitely many primes.
\end{corollary}

\begin{proof}
Take $D=\ell$ in Theorem~\ref{thm:structural}.
\end{proof}

\begin{remark}
If every prime divisor of $A_q$ is $1\pmod3$, Corollary~\ref{cor:pf} does not apply.  This is only an obstruction to the present safe-residue construction; it is not claimed to prevent exact-depth realization by another mechanism.
\end{remark}

\section{Density one of the structural family}
We now show that the hypothesis of Corollary~\ref{cor:pf} holds for almost all prime parameters $q\equiv1\pmod4$.

\begin{theorem}[density one]\label{thm:density}
Among primes $q\equiv1\pmod4$, the set for which
\[
A_q=\frac{3q+1}{4}
\]
has a prime divisor congruent to $2\pmod3$ has relative natural density $1$.
\end{theorem}

\begin{proof}
Avoiding the prime factor $2$ means $A_q$ is odd.  Among primes $q\equiv1\pmod4$, the two reduced classes modulo $8$ are $1$ and $5$, and
\[
2\mid A_q\iff q\equiv5\pmod8.
\]
Thus the relative density contribution from avoiding $2$ is $1/2$.

Now fix an odd prime $\ell\equiv2\pmod3$.  Since $\ell\ne3$,
\[
\ell\mid A_q
\iff
q\equiv-3^{-1}\pmod\ell. \tag{6}
\]
The residue in (6) is nonzero.  Conditional on $q\equiv1\pmod4$, the prime number theorem in arithmetic progressions modulo $4\ell$ shows that avoiding this one residue modulo $\ell$ has relative density
\[
1-\frac1{\ell-1}.
\]

For any finite set $S$ of odd primes $\ell\equiv2\pmod3$, CRT and the prime number theorem in arithmetic progressions give the exact limiting relative density
\[
\frac12\prod_{\ell\in S}\left(1-\frac1{\ell-1}\right) \tag{7}
\]
for primes $q\equiv1\pmod4$ whose $A_q$ is divisible by none of $2$ or the primes in $S$.

The classical divergence
\[
\sum_{\ell\equiv2(3)}\frac1\ell=\infty
\]
implies
\[
\prod_{\ell\in S}\left(1-\frac1{\ell-1}\right)\to0
\]
as $S$ exhausts the primes $2\pmod3$.  Hence the upper relative density of the exceptional set is zero.
\end{proof}

\begin{corollary}[density-one composite backbone]
For relative density $1$ of primes $q\equiv1\pmod4$, the composite depth
\[
\boxed{A_q=\frac{3q+1}{4}}
\]
is realized as exact square-completion depth by infinitely many primes.
\end{corollary}

\section{Quantitative consequence}
Define
\[
N_{\rm comp}(X)=\#\{A\le X:4A-1\text{ composite and }\delta_{\rm sq}(p)=A\text{ for infinitely many primes }p\}.
\]

\begin{theorem}[counting lower bound]\label{thm:count}
\[
\boxed{
N_{\rm comp}(X)\ge\left(\frac23+o(1)\right)\frac{X}{\log X}.
}
\]
\end{theorem}

\begin{proof}
The map $q\mapsto(3q+1)/4$ is injective.  By Theorem~\ref{thm:density} and the prime number theorem in the class $1\pmod4$,
\[
\#\{q\le Y:q\text{ prime},\ q\equiv1\pmod4,\ A_q\in\mathcal D^{\rm comp}_\infty\}
=\left(\frac12+o(1)\right)\frac{Y}{\log Y}.
\]
Now $A_q\le X$ if and only if $q\le(4X-1)/3$.  Substitution gives the result.
\end{proof}

\section{Explicit fixed-witness subfamilies}
For completeness we record a convenient explicit parametrization.

\begin{corollary}
Fix $d\equiv2\pmod3$ with $3\nmid d$.  Let $a_d\pmod{4d}$ be the unique reduced class satisfying
\[
3a_d\equiv-1\pmod{4d}.
\]
Every prime $q\equiv a_d\pmod{4d}$ gives an exact depth
\[
A_q=\frac{3q+1}{4}
\]
with witness $D=d$, realized by infinitely many primes.
\end{corollary}

The simplest case is $d=2$, which gives $q\equiv5\pmod8$.  For $q>5$ in this subfamily, the exact-depth prime progression may be chosen inside
\[
p\equiv1\pmod{840}.
\]

\section{General finite CRT certificates}
The special target modulus $3q$ allows the symbolic proof above.  For a general composite target one can still reduce exact first-hit to a finite congruence problem.

Fix $A,D$ with $D\mid A^2$, put $m_A=4A-1$ and $a\equiv-4D\pmod{m_A}$, and write
\[
n=a+m_Ak. \tag{8}
\]
An earlier witness $(B,E)$ imposes
\[
n\equiv-4E\pmod{m_B},
\qquad m_B=4B-1. \tag{9}
\]
Let $g=\gcd(m_A,m_B)$.  The progressions (8)--(9) intersect if and only if
\[
a\equiv-4E\pmod g.
\]
When they do, exactly one residue class of $k$ is removed modulo $m_B/g$.

\begin{proposition}[finite exact-depth certificate]
If the union of all earlier forbidden $k$-classes leaves a residue class $k\equiv r\pmod L$ such that
\[
\gcd(a+m_Ar,m_AL)=1,
\]
then infinitely many primes have exact square-completion depth $A$.
\end{proposition}

\begin{proof}
The resulting reduced progression contains the target witness and no earlier witness; Dirichlet's theorem supplies infinitely many prime members.
\end{proof}

The accompanying code represents $r$ by maximal prime-power coordinates rather than by one enormous CRT integer and verifies all earlier constraints from first principles.

\section{Certified composite full layers}
Two examples illustrate behavior beyond the symbolic $3q$ family.

At
\[
A=649,
\qquad
4A-1=2595=3\cdot5\cdot173,
\]
the realized divisor-modulus primitive branches
\[
D\in\{11,59,7139,38291\}
\]
all admit reduced exact-depth certificates.  Thus every one supports infinitely many primes with $\delta_{\rm sq}(p)=649$.

At
\[
A=22339,
\qquad
4A-1=89355=3\cdot5\cdot7\cdot23\cdot37,
\]
the four realized primitive branches
\[
D\in\{89,251,22361,67079\}
\]
likewise all support infinitely many exact-depth prime families.

These computations are supplementary.  The structural and density-one theorems above are computation-free.

\section{Relation to prior work and open problems}
The divisor-square Type-II mechanism itself is prior art, with historical provenance through Th\'epault and later expositions by Mizony and Gardes.  Modern complete Type-II parametrizations include work of Elsholtz--Tao and Pomerance--Weingartner.  Dahan studies Type-II sieve dimension and search depth.

Public 2026 FCF/CENTL work develops a prime-modulus exact-depth backbone, square-completion shadow theory, structural gaps, and a composite residual core.  The present candidate contribution is the composite $3q$ backbone, its prime-factor criterion, the density-one theorem, and the counting bound of Theorem~\ref{thm:count}.  Targeted searches through August 30, 2026 did not locate these exact statements; that negative search does not establish publication priority.

Natural next questions include:
\begin{enumerate}
\item Obtain a sharper asymptotic for $N_{\rm comp}(X)$, beyond the present lower bound.
\item Classify the exceptional prime parameters $q\equiv1\pmod4$ for which every prime factor of $(3q+1)/4$ is $1\pmod3$.
\item Find an analogue with target modulus $rq$ for $r>3$.  The present proof explains why $3$ is special: for $r>3$, the factor $q$ can already occur in the earlier modulus $3q$.
\item Determine whether infinitely many composite layers have complete realized primitive-layer realization, as in the certified examples $649$ and $22339$.
\end{enumerate}

The results here do not imply universal Type-II solvability.  The Erd\H{o}s--Straus conjecture remains open.

\begin{thebibliography}{9}
\bibitem{ET}
C. Elsholtz and T. Tao,
Counting the number of solutions to the Erd\H{o}s--Straus equation on unit fractions,
\emph{J. Aust. Math. Soc.} 94 (2013), 50--105.

\bibitem{PW}
C. Pomerance and A. Weingartner,
Exceptions to the Erd\H{o}s--Straus--Schinzel conjecture,
\emph{Ramanujan J.} 69 (2026).

\bibitem{Dahan}
B. Dahan,
Sieve dimension and search depth for the Erd\H{o}s--Straus conjecture, $n\equiv1\pmod{24}$,
arXiv:2608.24035 (2026).

\bibitem{MG}
L. Mizony and J.-L. Gardes,
Un point sur la conjecture d'Erd\H{o}s et Straus,
Universit\'e Claude-Bernard Lyon 1 (2012).

\bibitem{FCF}
Free Computation Foundation / CENTL,
Erd\H{o}s--Straus research library: square completion, exact-depth spectrum, and current frontier (2026).
\end{thebibliography}

\end{document}
